\chapter{Standard internazionali e framework per la cybersicurezza}
\label{cap:standard}

% 2.1 %%%%%%%%%%%%%%%%%%%%%%%%%%%%%%%%%%%%%%%%%%%%%%%%%%%%%%%%%%%%%%%%%%%%%%%%
\section{Standard ISO/IEC 27001:2022}

\subsection{Scopi e struttura dello standard}
L’ISO/IEC 27001:2022 è lo standard internazionale di riferimento per la gestione della sicurezza delle informazioni,
con oltre 40.000 certificazioni rilasciate a livello mondiale\footcite[p. 2]{mcintosh}.

Definisce i requisiti per l’istituzione, l’attuazione, il mantenimento e il miglioramento continuo di un 
\textit{Sistema di Gestione per la Sicurezza delle Informazioni} (SGSI), con l’obiettivo di garantire la protezione della riservatezza, 
integrità e disponibilità delle informazioni trattate dalle organizzazioni\footcite[pp. iv--v]{iso27001}. 
Lo standard si propone non soltanto come mezzo tecnico, ma anche come strumento organizzativo per strutturare la sicurezza 
informatica nei processi e negli aspetti di governance.

Lo standard segue \textit{Annex SL}, ovvero una struttura di alto livello comune a tutti gli standard ISO per i sistemi di gestione.
Tale struttura prevede dieci capitoli\footcite{uni-harmonized-structure}: 
\begin{enumerate}
    \item Scopo
    \item Riferimenti normativi
    \item Termini e definizioni
    \item Contesto dell'organizzazione
    \item Leadership
    \item Pianificazione
    \item Supporto
    \item Attività operative
    \item Valutazione delle prestazioni
    \item Miglioramento
\end{enumerate}
Inoltre, si fonda sul principio del miglioramento continuo, attuato attraverso il ciclo PDCA 
(Plan, Do, Check, Act), che guida in modo sistematico la pianificazione, l’implementazione, 
la verifica e il perfezionamento dei processi di gestione della sicurezza delle informazioni\footcite[p. 17]{watkins-iso}.

A corredo del contenuto dello standard, l'\textit{Annex A} riporta un catalogo di controlli 
di sicurezza selezionabili dalle aziende in base al proprio profilo di rischio\footcite[Annex A]{iso27001},
adattando così le misure alle necessità di ogni SGSI.


\subsection{Misure di gestione del rischio informatico}
La gestione del rischio costituisce l’elemento fondante di ISO 27001. Lo standard richiede infatti che l’organizzazione 
sia dotata di un processo strutturato per identificare, valutare e trattare i rischi per la sicurezza delle informazioni, in coerenza 
con il proprio contesto e con gli obiettivi strategici\footcite[pp. 6--8]{iso27001}.

Il processo di gestione del rischio si articola in due fasi principali. La prima è la valutazione del rischio (\textit{risk assessment}), 
che comprende l’identificazione delle minacce e delle vulnerabilità, l’analisi delle probabilità e degli impatti di potenziali incidenti
ai danni dell'azienda e la determinazione dei livelli di rischio accettabili.
La seconda è il trattamento del rischio (\textit{risk treatment}), in cui l’organizzazione definisce e 
implementa misure tecniche e procedurali per ridurre i rischi a un livello ritenuto tollerabile\footcite[pp. 4--5]{iso27001}.
Tale livello non è stabilito dallo standard, ma deve essere determinato dalle singole imprese in relazione alla loro esposizione
alle minacce e alla capacità di tollerare il rischio, secondo il risk framework da loro adottato.

A completamento del processo, lo standard stabilisce che i risultati della gestione del rischio debbano essere documentati e 
integrati nei processi di governance e di pianificazione strategica. La responsabilità ultima della valutazione e accettazione dei 
rischi spetta alla direzione, la quale deve dimostrare un impegno attivo nella supervisione del sistema di gestione della sicurezza 
delle informazioni\footcite[pp. 5--6]{iso27001}.


\subsection{Controlli di sicurezza: Annex A}
L’\textit{Annex A} rappresenta il principale strumento operativo a supporto del trattamento del rischio:
raccoglie un insieme di controlli di sicurezza che le organizzazioni possono adottare come misure di trattamento dei 
rischi individuati nei processi di valutazione\footcite[pp. iv--v]{iso27001}.
Ogni controllo è formulato come un obiettivo da perseguire, lasciando all’organizzazione la libertà di definire 
le modalità più adeguate di implementazione in base al proprio contesto\footcite[p. 16]{watkins-iso}.

I controlli elencati sono 93 e sono organizzati in quattro domini tematici principali\footcite[pp. 28--29]{watkins-iso}:
\begin{enumerate}
    \item Controlli organizzativi: includono politiche di sicurezza, responsabilità, relazioni con fornitori, gestione dei progetti e continuità operativa.
    \item Controlli per le persone: comprendono la consapevolezza del personale, la formazione e la definizione dei ruoli di sicurezza.
	\item Controlli fisici: riguardano la protezione degli asset fisici, delle infrastrutture e degli ambienti di lavoro.
	\item Controlli tecnologici: coprono ambiti come la gestione degli accessi, la protezione dei dati, il monitoraggio e la crittografia.
\end{enumerate}

L'\textit{Annex A}, pur non essendo esaustivo o prescrittivo, rappresenta un punto di riferimento comune per la 
progettazione delle misure di sicurezza. La sua funzione è quella di fornire un repertorio di \textit{best practices} utilizzate e riconosciute a 
livello internazionale che possono essere selezionate, adattate o integrate in funzione delle necessità specifiche dell'impresa. 
Questo elemento rappresenta uno dei punti di forza di ISO, poiché permette alle imprese di calibrare le proprie 
politiche di sicurezza sulla base del profilo di rischio specifico del SGSI\footcite[p. 22]{watkins-iso}.

I controlli diventano anche parte integrante dello \textit{Statement of Applicability} (SoA), il principale
documento obbligatorio per la certificazione ISO che esplicita le misure adottate e ne giustifica l’esclusione ove non pertinenti\footcite[pp. 4--5]{iso27001}.


\subsection{Implicazioni organizzative della certificazione ISO 27001}
L’adozione di ISO 27001 incide profondamente sulla governance, 
sui processi aziendali e sull'organizzazione in generale. Lo standard richiede il coinvolgimento attivo della direzione nella definizione della 
politica di sicurezza, nell’allocazione delle risorse e nella supervisione del SGSI, in conformità con le clausole 5.1 e 5.2\footcite[pp. 2--3]{iso27001}. 
L’adesione alla struttura armonizzata \textit{Annex SL} rafforza questa centralità, rendendo la sicurezza parte integrante del governo aziendale e facilitandone 
l’integrazione con gli altri sistemi di gestione\footcite{uni-harmonized-structure}.

Dal punto di vista dei processi, la certificazione implica l'adozione di meccanismi di valutazione delle prestazioni, audit e riesame della direzione (clausole 9.1--9.3), 
che costituiscono un feedback periodico sull'efficacia delle misure di sicurezza e i rischi da monitorare\footcite[pp. 8--9]{iso27001}\footcite[p. 2]{mcintosh}. 
Tali attività sono strumenti chiave per rafforzare anche l'accountability e il monitoraggio da parte del management.

La Dichiarazione di Applicabilità (SoA) costituisce un documento cruciale di trasparenza gestionale: collega i rischi individuati con i controlli adottati
(o esclusi, motivatamente), rendendo possibile al top management e agli auditor verificare la coerenza della strategia di sicurezza\footcite[pp. 4--5]{iso27001}.
In questo senso, l'\textit{Annex A} non rappresenta un insieme prescrittivo, ma un repertorio di misure da selezionare in base al profilo di rischio aziendale, 
responsabilizzando il top management nella definizione delle priorità per il SGSI da certificare\footcite[pp. 15--16]{watkins-iso}.

La certificazione impone anche requisiti in materia di competenze, consapevolezza e comunicazione (clausole 7.2--7.4), che richiedono programmi di formazione, 
definizione chiara dei ruoli e processi di comunicazione interna ed esterna sulla sicurezza\footcite[p. 6]{iso27001}. Questo favorisce lo 
sviluppo di una cultura aziendale della sicurezza, base della promozione di buone pratiche di igiene informatica.

Un ulteriore aspetto riguarda il miglioramento continuo (clausola 10), che formalizza la necessità di adattare politiche e controlli a nuove minacce 
e cambiamenti di contesto\footcite[p. 10]{iso27001}. Qui ISO sottolinea l'importanza di ripetere periodicamente i processi di
esame dei rischi e delle minacce, e di pianificazione delle misure necessarie, suggerendo quindi un ciclo di sicurezza aziendale, pur senza formalizzarlo.

ISO 27001 diventa così una leva manageriale per rendere tracciabili le decisioni sulla sicurezza, aumentare la fiducia degli stakeholder e sostenere la 
reputazione aziendale\footcite[p. 7]{watkins-iso}, oltre a fornire una metodologia strutturata per affrontare le minacce cyber. 
In prospettiva, la combinazione tra standardizzazione, audit e reporting facilita anche la dimostrazione di adempimento rispetto agli obblighi normativi, 
ben integrandosi con la Direttiva NIS2\footcite{enisa-implementing-guidance}\footcite{clusit-report}.



% 2.2 %%%%%%%%%%%%%%%%%%%%%%%%%%%%%%%%%%%%%%%%%%%%%%%%%%%%%%%%%%%%%%%%%%%%%%%%
\section{NIST Cybersecurity Framework 2.0}

\subsection{Scopi e struttura del framework}
Il \textit{NIST Cybersecurity Framework} (CSF), a cura del \textit{National Institute of Standards and Technology} (NIST), nasce nel 2014
con l’obiettivo di migliorare la protezione delle infrastrutture critiche negli Stati Uniti. La nuova versione 2.0, 
pubblicata nel febbraio 2024, amplia la portata del framework, rendendolo uno strumento utilizzabile a livello 
internazionale e per organizzazioni di ogni dimensione e settore\footcite[p. 1]{nist-csf}. Il NIST CSF, come l’ISO 27001, 
è uno dei framework di cybersicurezza più diffusi a livello globale, grazie alla sua capacità di adattamento a diversi settori 
e alla funzione di ponte tra aspetti tecnici e manageriali\footcite[p. 1]{mcintosh}.

La struttura del CSF si articola in tre componenti principali:
\begin{itemize}
    \item \textbf{Core}, il nucleo del framework, che raccoglie un insieme di \textit{outcome} organizzati in Funzioni, 
    Categorie e Sottocategorie. In particolare, sono individuate sei funzioni principali: 
    \textit{Govern}, \textit{Identify}, \textit{Protect}, \textit{Detect}, \textit{Respond}, \textit{Recover}\footcite[p. 5]{nist-csf} (Fig. \ref{fig:funzioni-nist}). 
    \item \textbf{Organizational Profiles}, che descrivono lo stato attuale e quello target del posizionamento di sicurezza 
    di un’organizzazione, consentendo di confrontare i gap e pianificare il miglioramento\footcite[p. 6]{nist-csf}.  
    \item \textbf{Tiers}, che rappresentano i livelli di maturità nella gestione del rischio, utili per 
    valutare la capacità dell’organizzazione di integrare la cybersecurity nei processi decisionali e culturali\footcite[p. 7]{nist-csf}.  
\end{itemize}

\defaultimage{nist-fig2}{Funzioni del framework NIST CSF 2.0\footcite[Fig. 2]{nist-csf}}{fig:funzioni-nist}

Lo scopo principale del CSF è fornire un modello flessibile che consenta alle imprese di 
gestire i rischi informatici in modo sistematico e integrato con gli altri rischi caratteristici del business\footcite[p. 2]{nist-csf}.
Un'organizzazione può infatti utilizzare gli elementi Core, Profiles e Tiers per\footcite[pp. 1--2]{nist-csf}:
\begin{itemize}
    \item descrivere la postura di cybersicurezza attuale o target di una parte o dell’intera organizzazione, determinare i gap e 
    valutare i progressi nel loro superamento;
    \item identificare, organizzare e dare priorità alle azioni per la gestione dei rischi di cybersicurezza in modo coerente con la missione 
    dell’organizzazione, i requisiti legali e regolamentari, nonché le aspettative di gestione del rischio e di governance;
	\item fornire un linguaggio comune per comunicare, all’interno e all’esterno dell’organizzazione, in merito a rischi di 
    cybersicurezza, capacità, esigenze ed aspettative.
\end{itemize}
Inoltre, il CSF 2.0 rafforza l’attenzione alla supply chain, riconosciuta come uno dei vettori di rischio più critici 
nello scenario contemporaneo\footcite{nist-csf}, suggerendo una metodologia strutturata per la comunicazione in
ambito di \textit{Third Party Risk Management}.

Dunque, a differenza di ISO, che stabilisce requisiti certificabili per quanto concerne la creazione, la 
documentazione e il mantenimento di un Sistema di Gestione per la Sicurezza delle Informazioni, il 
NIST CSF adotta un approccio più ampio, definendo misure di sicurezza sulla base di \textit{outcome} 
desiderati\footcite[pp. 3--4]{nist-csf} per quanto concerne tutti gli aspetti della cybersicurezza aziendale.


\subsection{Approccio basato sul rischio e implementazione}
La gestione del rischio si articola attraverso il CSF Core, che raccoglie sei funzioni principali: 
\textit{Govern}, aggiunta nella versione 2.0 del CSF,
\textit{Identify}, \textit{Protect}, \textit{Detect}, \textit{Respond} e \textit{Recover}\footcite[p. 5]{nist-csf} (Fig. \ref{fig:funzioni-nist}). 
Esse sono organizzate in categorie e sottocategorie, contenute nell'Appendice A del framework,
che traducono gli obiettivi generali delle funzioni in risultati precisi e misurabili, che le aziende
implementano mediante controlli di sicurezza.
\begin{itemize}
    \item \textit{Govern} enfatizza la necessità di collegare la cybersecurity alla governance aziendale, 
    integrando la gestione del rischio informatico con la strategia complessiva dell’impresa\footcite[p. 8]{nist-csf}.
    \item \textit{Identify} si concentra sulla conoscenza del contesto organizzativo, degli asset critici e delle 
    dipendenze, al fine di valutare correttamente le minacce.
    \item \textit{Protect} è invece orientata all’implementazione di misure preventive, 
    come il controllo degli accessi, la protezione dei dati e la formazione del personale.
    \item \textit{Detect} ha lo scopo di sviluppare capacità di monitoraggio e di individuazione tempestiva 
    di eventi e anomalie di sicurezza.
    \item \textit{Respond} definisce le modalità di gestione degli incidenti, dalla comunicazione alla mitigazione.
    \item \textit{Recover} garantisce il ripristino delle capacità operative e l’attivazione di meccanismi di 
    miglioramento post-incidente, rafforzando la resilienza organizzativa.
\end{itemize}

Accanto al Core, il framework prevede due strumenti fondamentali di implementazione: i 
\textit{Profiles} e gli \textit{Implementation Tiers}. I \textit{Profiles} descrivono lo stato attuale e lo stato 
target della postura di sicurezza di un’organizzazione, consentendo di individuare le carenze nei confronti dell'obiettivo desiderato e 
pianificare un percorso di miglioramento\footcite[p. 6]{nist-csf}. I \textit{Tiers}, articolati in quattro livelli 
(dal \textit{Partial} all’\textit{Adaptive}), rappresentano invece la maturità dei processi di gestione del rischio 
e la capacità di integrazione della cybersecurity nei processi decisionali e culturali. Pur non essendo 
certificativi, rappresentano validi strumenti di autovalutazione per l'azienda\footcite[p. 7]{nist-csf}. 

L’implementazione pratica del CSF attraversa ciclicamente cinque fasi: 
\begin{enumerate}
    \item comprensione del contesto e della missione aziendale;
    \item definizione dei profili e analisi dei gap;
    \item prioritizzazione degli interventi sulla base del rischio e delle risorse disponibili;
    \item applicazione delle misure individuate;
    \item monitoraggio e revisione continua.
\end{enumerate}

Dunque, come avviene con il modello PDCA tipico degli standard ISO, il processo mantiene una logica di miglioramento 
continuo, volta ad adattare le misure a minacce e tecnologie in evoluzione.

Il \textit{NIST Cybersecurity Framework}, nonostante la non certificabilità a differenza di ISO,
è oggi ampiamente diffuso a livello globale grazie alla sua 
capacità di fungere da punto di raccordo tra aspetti tecnici e manageriali\footcite[p. 3]{mcintosh}. 
Un altro punto di forza risiede nella flessibilità: i vari documenti di supporto pubblicati dal NIST per 
l’implementazione del framework vengono aggiornati con regolarità, consentendo di adeguare le misure 
di sicurezza all’evoluzione delle minacce e del contesto tecnologico.
Parmar e Miles ne sottolineano la crescente rilevanza in ambito europeo, dove è sempre 
più utilizzato come riferimento di \textit{best practice} e come base per la mappatura con la Direttiva 
NIS2\footcite[p. 12]{parmar-miles-csf}.


\subsection{Implicazioni organizzative dell'adozione del NIST CSF 2.0}
Dal punto di vista organizzativo, il NIST CSF definisce un processo di gestione e trattamento del rischio
ciclico. Matteo Herin osserva come le funzioni del framework siano rappresentate secondo un cerchio, alla cui base si colloca la 
funzione di governo e intorno alla quale si susseguono le altre. Il primo step è rappresentato 
dall'identificazione delle minacce, per poi procedere con l'implementazione delle protezioni necessarie.
Le funzioni successive entrano in gioco in caso di attacco, con il monitoraggio attivo e
la risposta a eventuali attacchi tramite i playbook codificati per contenere ed estirpare l'ostilità.
Infine, avviene il recupero dall'incidente con il ripristino delle normali funzioni, per poi ricominciare
il ciclo dalla funzione \textit{Identify} sulla base della \textit{lesson learned} ricavata
dall'analisi degli eventi.

Un'attenzione particolare deve essere rivolta alla funzione \textit{Govern}, i cui obiettivi risultano 
coerenti con le priorità individuate dagli enti europei\footcite[Allegato B]{enisa-threat-landscape}.
Essa evidenzia la necessità di un coinvolgimento esplicito della 
direzione nella definizione delle strategie di sicurezza, nella distribuzione 
delle responsabilità e nella supervisione dei processi di gestione del rischio\footcite[Appendice A]{nist-csf}. 
In questo senso, il CSF si propone come strumento per rafforzare il legame tra obiettivi 
strategici e misure operative, rendendo la sicurezza informatica parte integrante della strategia aziendale.

L’utilizzo dei \textit{Profiles} e dei \textit{Tiers} consente al management 
di valutare lo stato di maturità della propria postura di sicurezza e di stabilire priorità di 
investimento. Ciò fornisce ai CISO uno strumento di dialogo con il vertice 
aziendale, facilitando la traduzione dei rischi tecnici in rischi di business e la pianificazione di 
roadmap di progresso\footcite{parmar-miles-csf}. Il framework diventa così anche un 
mezzo di comunicazione strategica, utile a colmare il divario tra linguaggio tecnico e 
decisioni manageriali, rendendo più efficace e diretta la comunicazione in materia di cybersicurezza.

L’adozione del CSF stimola inoltre la diffusione di una cultura organizzativa del rischio, 
favorendo la sensibilizzazione del personale, l’assegnazione di ruoli chiari e la creazione di 
processi di feedback continui. La sua implementazione rafforza la resilienza aziendale e contribuisce alla dimostrazione di 
\textit{due diligence}, aspetto particolarmente rilevante in un contesto regolatorio come quello 
della Direttiva NIS2, che richiede alle organizzazioni di integrare la cybersecurity nella propria 
struttura di governance\footcite[pp. 110--115]{enisa-threat-landscape}. 
Il framework si configura dunque come un supporto essenziale per l’IT Manager, chiamato a guidare 
l'implementazione delle linee guida NIST attraverso misure concrete.
